\documentclass[10pt,a4paper]{amsart}
\usepackage[margin=19mm]{geometry}
\usepackage{amsmath,amssymb,mathtools}
\usepackage[scr=rsfs,cal=euler]{mathalfa}
\usepackage{xfrac}
\usepackage[unicode,hidelinks,bookmarks=false]{hyperref}
\newcommand{\ie}{i.e.\ }
\newcommand{\cf}{cf.\ }
\newcommand{\resp}{respectively\ }
\newcommand{\Subsection}{\subsection}
\providecommand{\nobreakdash}{\nobreak\hbox{-}}
\setlength{\emergencystretch}{2em}
\multlinegap0pt
\usepackage{xcolor}
\usepackage{braket}
\usepackage{amssymb}
\usepackage{mathtools}
\usepackage[shortlabels]{enumitem}
\usepackage{bm}
\usepackage{float}
\newtheorem{proposition}{Proposition}
\newcommand{\cB}{\mathcal B}
\title{Approximate QCAs in one dimension using approximate~algebras}
\author{Daniel Ranard, Michael Walter, Freek Witteveen}
\date{Source: arXiv:2603.08702v1 (CC BY 4.0)}
\begin{document}
\maketitle

\noindent\emph{Source and licence.} Approximate QCAs in one dimension using approximate~algebras. Version arXiv:2603.08702v1 (CC BY 4.0), distributed by arXiv under the Creative Commons Attribution 4.0 International licence, http://creativecommons.org/licenses/by/4.0/, as recorded in the arXiv licence metadata for this exact version and shown on the abstract page https://arxiv.org/abs/2603.08702v1. Full licence notice: \texttt{licenses/arxiv-2603.08702-CC-BY-4.0.txt}.\par\smallskip

\noindent\textit{Editorial scope.} Counted unit: the article's proposition prop:almost-unitary-polar (source0.tex lines 1338--1351). The article's complete original proof of it is delivered with it (source0.tex lines 1353--1396, one \texttt{proof} environment of 44 lines). No mathematics is written, repaired or re-derived: the statement and the proof are the article's own lines, in the article's own notation, and no retained line is substituted or re-flowed.  No further source-local context is needed: every cross-reference of every retained line resolves inside the two delivered spans.  Nothing was omitted from any retained span, so the package carries no omission record.  No retained line contains a citation, so the wrapper carries no bibliography and no bibliography entry.  No retained line contains a figure environment, an image-inclusion command or drawing code, so no image is delivered and the package declares no asset dependency.  Editorial formatting only, in addition: an AMS \texttt{amsart} preamble that restates the article's own declarations and macros used by the retained lines, namely the environments \texttt{proposition} on separate counters, together with the standard packages those lines need.  The retained environments print in this document as Proposition 1 in order of appearance, whereas the article prints the same items with the numbering of its own full document.  The complete native source of the article as distributed in the arXiv e-print for this exact version is bundled unchanged as \texttt{sources/196025/source0.tex}.
\begin{proposition}[Almost-unitaries are close to unitaries]\label{prop:almost-unitary-polar}
Let $\cB$ be a unital finite-dimensional $C^*$-algebra, and let $x \in \cB$ satisfy
\[
\|x^*x-I\| \le \epsilon < \frac12.
\]
Then $x^*x$ is invertible, so the polar part
\[
u := x(x^*x)^{-1/2}
\]
is a unitary in $\cB$, and
\begin{align}
\|x-u\| \le 2\epsilon.
\end{align}
\end{proposition}

\begin{proof}
Since $\|x^*x-I\|<1$, the spectrum of $x^*x$ is contained in $[1-\epsilon,1+\epsilon] \subset (0,\infty)$, so $x^*x$ is invertible. Hence $u=x(x^*x)^{-1/2} \in \cB$ is well-defined and satisfies
\[
u^*u=(x^*x)^{-1/2}x^*x(x^*x)^{-1/2}=I.
\]
Since $x^*x$ is invertible, also $xx^*$ is invertible, so $u$ is unitary. Note
\[
\|x-u\|
=
\|x-x(x^*x)^{-1/2}\|
\le
\|x\|\,\|I-(x^*x)^{-1/2}\|.
\]
Also,
\[
\|x\|^2=\|x^*x\| \le 1+\epsilon,
\]
so $\|x\| \le (1+\epsilon)^{1/2}$. Finally, since $\sigma(x^*x)\subset [1-\epsilon,1+\epsilon]$,
\[
\|(x^*x)^{-1/2}-I\|
\le
\max_{t\in[1-\epsilon,1+\epsilon]} |t^{-1/2}-1|
=
(1-\epsilon)^{-1/2}-1.
\]
Therefore
\[
\|x-u\|
\le
(1+\epsilon)^{1/2}\left((1-\epsilon)^{-1/2}-1\right)
=
\frac{(1+\epsilon)^{1/2}-(1-\epsilon)^{1/2}}{(1-\epsilon)^{1/2}}.
\]
With $\epsilon<1/2$,  we obtain the desired $\|x-u\|\le 2\epsilon$.
% we have $(1-\epsilon)^{-1/2}\le \sqrt{2}$, and
% \[
% (1+\epsilon)^{1/2}-(1-\epsilon)^{1/2}
% =
% \frac{2\epsilon}{(1+\epsilon)^{1/2}+(1-\epsilon)^{1/2}}
% \le
% \epsilon.
% \]
% Hence we obtain the desired $\|x-u\|\le 2\epsilon.$
\end{proof}

\end{document}
